% !TeX root = ../../Book4.tex
% 纲要标题：### 14.4 正合偶与低阶计算
% 纲要位置：计划纲要.md，第 2845 行。

\section{正合偶与低阶计算}
\label{b4:sec:1404}

前两节已从滤过复形直接构造谱序列并证明有限收敛。本节用正合偶重新组织同一过程，再考察非零行列很少时的微分、边缘同态和扩张问题。

% 纲要内容：
%
% * 正合偶的构造
% * 两列、三列和第一象限例子
% * 退化与边缘同态
% * 不忽略隐藏延拓问题
%

\subsection{正合偶的构造}
\label{b4:subsec:1402:02}

谱序列还有一种以长正合列为起点的构造。它把“微分后再看商”的过程压缩成三个相互正合的箭头。

\begin{definition}\label{b4:def:exact-couple}
在交换范畴中，给定对象 \(D,E\) 与态射 \(i:D\to D\)、\(j:D\to E\)、\(k:E\to D\)。若
\[
 \operatorname{im}i=\ker j,\quad
 \operatorname{im}j=\ker k,\quad
 \operatorname{im}k=\ker i,
\]
则称这些数据构成\term[zhenghe-ou]{正合偶}[exact couple]。分次情形允许三个态射具有指定的次数。
\end{definition}

\begin{proposition}\label{b4:prop:derived-exact-couple}
设 \(R\) 为环，\((D,E,i,j,k)\) 是左 \(R\)-模范畴中的正合偶，\(i:D\to D\)、\(j:D\to E\)、\(k:E\to D\)。令 \(\mathrm{d}=jk\)，则 \(\mathrm{d}^2=0\)。取
\[
 D'=\operatorname{im}i,\quad E'=\ker \mathrm{d}/\operatorname{im}\mathrm{d},
\quad i'=i|_{D'},\quad j'(ia)=[ja],\quad k'([e])=ke.
\]
这些公式良定义，并给出新的正合偶。
\end{proposition}
\begin{proof}
因为 \(kj=0\)，有 \(\mathrm{d}^2=0\)。若 \(ia=ia'\)，则 \(a-a'=kb\)，所以 \(ja-ja'=\mathrm{d}b\)；又 \(\mathrm{d}ja=0\)，故 \(j'\) 良定义。若 \(\mathrm{d}e=0\)，则 \(ke\in\ker j=\operatorname{im}i\)；若更改 \(e\) 一个 \(\mathrm{d}b\)，则 \(k\mathrm{d}b=kjkb=0\)，所以 \(k'\) 良定义。

若 \(j'(ia)=0\)，则 \(ja=jkb\)，于是 \(a-kb=ic\)；因此 \(ia=i^2c\)，这证明 \(\ker j'=\operatorname{im}i'\)。若 \(k'[e]=0\)，则 \(ke=0\)，可写 \(e=ja\)，从而 \([e]=j'(ia)\)，给出第二处正合。最后，若 \(ia\in D'\) 且 \(i(ia)=0\)，写 \(ia=ke\)。由于 \(jke=jia=0\)，\(e\) 是 \(\mathrm{d}\)-余圈，且 \(k'[e]=ia\)。各相邻复合显然为零，所以第三处也正合。以上过程只涉及核、像和有限提升，也证明交换范畴版本。
\end{proof}

应用于短正合复形列 \(0\to F^{p+1}C\to F^pC\to\operatorname{gr}_F^pC\to0\) 的长正合列，得到
\[
 D_1^{p,q}=H^{p+q}(F^pC),\qquad E_1^{p,q}=H^{p+q}(\operatorname{gr}_F^pC).
\]
三个箭头依次为
\[
 i:D_1^{p+1,q-1}\to D_1^{p,q},\quad
 j:D_1^{p,q}\to E_1^{p,q},\quad
 k:E_1^{p,q}\to D_1^{p+1,q}.
\]
故 \(jk\) 正是 \(\mathrm{d}_1\)。反复导出正合偶时，每次把 \(D\) 替成包含映射的像；表示连接像所需的提升每次多经过一层滤过，得到次数 \((r,1-r)\) 的 \(\mathrm{d}_r\)。这些页与\cref{b4:eq:spectral-page-formula}一致：连接同态以 \(\mathrm{d}x\) 表示，而提升到第 \(r\) 层正要求 \(\mathrm{d}x\in F^{p+r}\)；更换提升产生的两类差正是该式分母。页间关系已经在\cref{b4:thm:filtered-spectral-sequence}中直接核对，因而正合偶记号不是省略计算的理由。

\subsection{两列、三列与第一象限例子}
\label{b4:subsec:1404:01}

若第一象限谱序列只在 \(p=0,1\) 两列非零，则从 \(r=2\) 起所有微分都没有可能的靶，故 \(E_2=E_\infty\)。对总次数 \(n\)，收敛给出
\[
 0\longrightarrow E_2^{1,n-1}\longrightarrow H^n
 \longrightarrow E_2^{0,n}\longrightarrow0.
\]
这说明两列计算天然产生一个扩张，并不自动产生两项直和。

\ProseParagraphBreak
若只在 \(p=0,1,2\) 三列非零，第一个仍可能出现的高阶微分是
\[
 \mathrm{d}_2:E_2^{0,q}\longrightarrow E_2^{2,q-1};
\]
它之后才有 \(E_3=E_\infty\)。\secref{b4:sec:1405}将把这个箭头真正算出来。若谱序列没有列数上界但在第一象限内，对固定 \((p,q)\) 的箭头仍会终止：当 \(r>q+1\) 时出箭头的靶在负行，当 \(r>p\) 时入箭头的源在负列。因此逐位置稳定，并不意味着整个谱序列在同一个有限页退化。

\subsection{退化与边缘同态}
\label{b4:subsec:1404:02}

\begin{definition}\label{b4:def:spectral-degeneration}
设 \(\mathcal A\) 为交换范畴，\(E_r^{p,q}\) 为 \(\mathcal A\) 中的上同调型谱序列，\(s\) 为其一个页数。若从页 \(s\) 开始的所有微分均为零，则称谱序列在第 \(s\) 页\term[tuihua-puxulie]{退化}[degeneration]。若谱序列位于第一象限并强收敛于带滤过 \(F\) 的 \(H^*\)，则对 \(n\ge0\)，目标滤过的端部诱导
\[
 E_2^{n,0}\longrightarrow E_\infty^{n,0}=F^nH^n\longrightarrow H^n,
 \qquad
 H^n\longrightarrow H^n/F^1H^n=E_\infty^{0,n}\longrightarrow E_2^{0,n},
\]
称为相应的\term[bianyuan-tongtai]{边缘同态}[edge homomorphisms]。
\end{definition}

这里第一条由各页的商映射组成，因为底行从第二页起只有入箭头；第二条由子对象包含组成，因为最左列只有出箭头。若谱序列从第零页给出，第一、二页仍可有竖直或水平微分，故这个端部描述特意从第二页开始。

\begin{proposition}\label{b4:prop:spectral-single-row}
设交换范畴中的第一象限上同调型谱序列强收敛于 \(H^*\)。若其第二页只在 \(q=0\) 行非零，则边缘同态 \(E_2^{n,0}\to H^n\) 对每个 \(n\ge0\) 为同构；若只在 \(p=0\) 列非零，则另一边缘同态 \(H^n\to E_2^{0,n}\) 对每个 \(n\ge0\) 为同构。
\end{proposition}
\begin{proof}
在第一种情形，第 \(r\ge2\) 页的出箭头靶行是 \(1-r<0\)，入箭头源在正行却为零，因此全部微分消失。目标滤过只有 \(\operatorname{gr}^nH^n\) 非零，有限性依次给出 \(F^nH^n=H^n\)、\(F^{n+1}H^n=0\)。第二种情形交换两坐标的角色，同样只剩一个非零相邻商。
\end{proof}

\subsection{隐藏延拓问题}
\label{b4:subsec:1404:03}

使用谱序列时可以把计算分成两件事：先求各页及其微分，得到 \(E_\infty\)；再利用目标的其他性质恢复滤过对象。后一步可能使用已知的生成元、乘法、指数或 \(\operatorname{Ext}^1\) 中的扩张类。它不属于再计算一个未出现的 \(\mathrm{d}_r\)。

\ProseParagraphBreak
例如已经知道有限群 \(H\) 的两个相邻商均为 \(\mathbb Z/2\)，则 \(|H|=4\) 已经确定；然而 \(H\) 的指数仍可能是二或四。若此外存在一个四阶元素，则 \(H\cong\mathbb Z/4\)；若 \(2H=0\)，则 \(H\cong(\mathbb Z/2)^2\)。判断原对象的指数，需要相邻商以外的扩张信息。
